%------------------------------------------------------------------------------%
\documentclass[fleqn,utf8,aspectratio=169,12pt]{beamer}

\usepackage{calculo_varias_variaveis-1}

\title{Funções de Várias Variáveis}

%------------------------------------------------------------------------------%
\begin{document}

\addtitle

%------------------------------------------------------------------------------%
\section{Funções}

%------------------------------------------------------------------------------%
\begin{frame}{Funções}
  \[
    f\colon A \to B
  \]
  \pause
  \vspace{-\baselineskip}
  \begin{itemize}[<+->]
    \item $A$ \emph{domínio}
    \item $B$ \emph{contradomínio}
    \item $f$ \emph{relação} que associa elementos do domínio a elementos
          do contradomínio
    \item Cada elemento do domínio deve estar associado a \emph{um único}
          elemento do contradomínio
    \item A \emph{imagem} é um subconjunto de $B$ composto dos
          \emph{elementos que são imagem} de algum elemento de $A$
  \end{itemize}
\end{frame}

%------------------------------------------------------------------------------%
\begin{frame}{Domínio}
  \begin{itemize}[<+->]
    \item Precisa ser dado a priori
    \item Geralmente é determinado pela aplicação
    \item Se nada for dito, assumimos o domínio ``natural'', \\
          ie, todos os valores onde podemos efetuar as contas
  \end{itemize}
\end{frame}

%------------------------------------------------------------------------------%
\begin{frame}{Imagem}
  \begin{itemize}[<+->]
    \item Subconjunto do contradomínio
    \item O contradomínio é definido a priori
    \item Pode ser difícil identificar a imagem de uma função
    \item Dada uma função \(f\colon D \to C \) \\
          $y$ pertence a imagem de $f$,
          se existe $x\in D$ tal que
          \[y=f(x)\]
  \end{itemize}
\end{frame}

%------------------------------------------------------------------------------%
\begin{frame}{Nomenclatura -- Curiosidade}
  Função real
  \[
    f\colon D \subset \R \to \R
  \]
  \pause
  Função de várias variáveis
  \[
    f\colon D \subset \R^n \to \R
  \]
  \pause
  Função vetorial ou paramétrica
  \[
    f\colon D \subset \R \to \R^n
  \]
  \pause
  Campo vetorial
  \[
    f\colon D \subset \R^n \to \R^m
  \]
\end{frame}

%------------------------------------------------------------------------------%
\section{Funções de Várias Variáveis}

%------------------------------------------------------------------------------%
\begin{frame}{Funções de Várias Variáveis}
  \[
    f\colon D \subset \R^n \to \R
  \]
  \pause
  \[
    f(x, y) = x + y
  \]
  \pause
  \[
    f(\vec{r}) = \sqrt{\sum_{k=1}^n r_k^2\,}
  \]
\end{frame}

%------------------------------------------------------------------------------%
\begin{frame}{Funções de Várias Variáveis}
  Suponha que $D$ seja um conjunto  de $n$-uplas de números reais
  \((x_1, x_2, \dots, x_n)\)

  \pause
  Uma \emph{função} a valores reais $f$ em $D$ é uma regra que associa
  \emph{um único número real}
  \[
    w = f(x_1, x_2, \dots, x_n)
  \]
  a cada elemento em $D$

  \pause
  O conjunto $D$ é o \emph{domínio} da função

  \pause
  O conjunto de valores de $w$ assumidos por $f$ é a \emph{imagem} da função

  \pause
  O símbolo $w$ é a variável \emph{dependente} de $f$, e dizemos que $f$
  é função de $n$ variáveis \emph{independentes} $x_1$ a $x_n$

  \pause
  Também chamamos os $x_j$ de variáveis de \emph{entrada} da função e
  denominamos $w$ a variável de \emph{saída} da função
\end{frame}

%------------------------------------------------------------------------------%
\begin{frame}{Domínio e Imagem}

  Determinamos o \emph{domínio} eliminando pontos que
  levam a divisões por zero ou valores complexos.

  \pause
  \bigskip

  A \emph{imagem} consiste on conjunto de valores de saída
  para a variável dependente
\end{frame}

%------------------------------------------------------------------------------%
\section{Exemplos}

%------------------------------------------------------------------------------%
%------------------------------------------------------------------------------%
\begin{question}
  Descreva o domínio da função \(f(x, y) = \sqrt{y-x^2}\)

  \bigskip
  \begin{enumerate}
    \item Qual a condição que os pontos do domínio deve satisfazer?
    \item Que região é essa? Represente-a graficamente.
  \end{enumerate}
\end{question}

%------------------------------------------------------------------------------%
\begin{resolution}{Solução}
  \onslide<+->{Condição que deve ser satisfeita para um ponto estar no domínio}
  \begin{align*}
    \onslide<+->{y-x^2 & \geq 0 \\}
    \onslide<+->{y & \geq x^2}
  \end{align*}

  \onslide<+->{
  Domínio
  \[
    D_f = \left\{ (x, y)\in\R^2 \st y \geq x^2\right\}
  \]}
\end{resolution}

%------------------------------------------------------------------------------%
\begin{resolution}{Representação Gráfica}
  \centering
  \includegraphics{ex_dominio_sqrt_y_x2}
\end{resolution}

%------------------------------------------------------------------------------%
%------------------------------------------------------------------------------%
\begin{question}
  Encontre e esboce o domínio da função
  \[
    f(x, y) = \ln\left(x^2 + y^2 - 4\right)
  \]
\end{question}

%------------------------------------------------------------------------------%
\begin{resolution}{Solução}
  \onslide<+->{Condição}
  \begin{align*}
    \onslide<+->{x^2 + y^2 - 4 & > 0 \\}
    \onslide<+->{x^2 + y^2 & > 4}
  \end{align*}

  \onslide<+->{Pontos exteriores ao circulo centrado na origem e raio 2}

  \onslide<+->{
  Domínio
  \[
    D_f = \left\{ (x, y)\in\R^2 \st x^2 + y^2 > 4\right\}
  \]}
\end{resolution}

%------------------------------------------------------------------------------%
\begin{resolution}{Representação Gráfica}
  \centering
  \includegraphics{ex_dominio_ln_x2_y2_4}
\end{resolution}

%------------------------------------------------------------------------------%
%------------------------------------------------------------------------------%
\begin{question}
  Encontre e esboce o domínio da função
  \(
    f(x, y) = \frac{\sqrt{x^2-4\,}}{\;\ln(y-1)\;}
  \)
\end{question}

%------------------------------------------------------------------------------%
\begin{resolution}{Solução}
  Condições para avaliar a função
  \(f(x, y) = \frac{\sqrt{x^2-4\,}}{\;\ln(y-1)\;}\)

  \pause
  \begin{enumerate}[<+->]
    \item \(x^2-4 \geq 0\)
    \item \(\ln(y-1) \neq 0\)
    \item \(y-1 > 0\)
  \end{enumerate}
\end{resolution}

%------------------------------------------------------------------------------%
\begin{resolution}{Solução}
\begin{columns}[t]
\column{0.3\linewidth}
  \onslide<+->{Condição 1}
  \begin{align*}
    \onslide<+->{x^2-4     & \geq 0                       \\}
    \onslide<+->{x^2       & \geq 4                       \\}
    \onslide<+->{\abs{x}   & \geq 2                       \\}
    \onslide<+->{x \leq -2 & \quad\text{ou}\quad x \geq 2}
  \end{align*}
\column{0.3\linewidth}
  \onslide<+->{Condição 2}
  \begin{align*}
    \onslide<+->{\ln(y-1) & \neq 0 \\}
    \onslide<+->{y-1      & \neq 1 \\}
    \onslide<+->{y        & \neq 2}
  \end{align*}
\column{0.3\linewidth}
  \onslide<+->{Condição 3}
  \begin{align*}
    \onslide<+->{y-1 & > 0 \\}
    \onslide<+->{y   & > 1}
  \end{align*}
\end{columns}
\end{resolution}

%------------------------------------------------------------------------------%
\begin{resolution}{Solução}
  Domínio
  \[
    D_f = \left\{ (x, y)\in\R^2 \st
      \abs{x} \geq 2 \;\text{ e }\;
      y > 1          \;\text{ e }\;
      y \neq 2
    \right\}
  \]
\end{resolution}

%------------------------------------------------------------------------------%
\begin{resolution}{Representação Gráfica}
  \centering
  \includegraphics{ex_dominio_sqrt_x2_4_ln_y_1}
\end{resolution}

%------------------------------------------------------------------------------%

%------------------------------------------------------------------------------%
\section{Lista Mínima}

%------------------------------------------------------------------------------%
\begin{frame}{Lista Mínima}
  Cálculo Vol. 2 do Thomas 12\textsuperscript{a} ed. -- Seção 14.1
  \begin{enumerate}
    \item Estudar o texto da seção
    \item Resolver os exercícios: 1-2, 5-9
  \end{enumerate}

  \vfill
  Atenção: A prova é baseada no livro, não nas apresentações
\end{frame}

%------------------------------------------------------------------------------%
\end{document}
%------------------------------------------------------------------------------%
